% !PRINTED 8
\weberpage{8}
%<<CONT
$Q_1$ are distinct. For if we had, say, $a_1 c = a_2 b$, it would follow that $c = a_1^{-1} a_2 b$; but $a_1^{-1} a_2$ is contained in $Q$, hence identical with one of the elements $1, a_1, a_2 \ldots a_{\nu-1}$, and $c$ would therefore, contrary to hypothesis, be contained in $Q_1$. We continue in this way, forming the systems $Q, Q_1, Q_2 \ldots Q_{\mu-1}$, of which the last shall be
\[
Q_{\mu-1} = g, a_1 g, a_2 g \ldots a_{\nu-1} g
\]
But then $P$ must be exhausted by the systems $Q, Q_1, Q_2 \ldots Q_{\mu-1}$, and since these systems contain nothing but distinct elements, it follows that
\[
n = \nu \mu,
\]
or the theorem:

\noindent 1.\ \gesperrt{The degree of a group is divisible by the degree of each of its divisors. The quotient $\mu = n : \nu$ shall be called the index of the divisor $Q$.}

The systems $Q_1, Q_2 \ldots Q_{\mu-1}$ we call, as in the special case of the permutation groups, the \gesperrt{subsidiary groups} belonging to $Q$, and denote them by
\[
Q_1 = Qb, \quad Q_2 = Qc \ldots Q_{\mu-1} = Qg,
\]
so that we can also write symbolically
\weq{1}{P = Q + Qb + Qc + \cdots + Qg}
At times we shall also designate the whole system $Q, Q_1 \ldots Q_{\mu-1}$ ($Q$ itself included) as a system of subsidiary groups, and accordingly call the representation (1) the \gesperrt{decomposition of $P$ into a system of subsidiary groups}.

From the mode of formation of the subsidiary groups it follows further that, if $b$, $c$ are any two elements from $P$, the two systems $Qb$ and $Qc$ are either entirely identical or contain no common element. For if $a_1 b = a_2 c$, where $a_1$, $a_2$ are two elements from $Q$, then for every other element $a$ from $Q$ we also have
\[
a a_1 b = a a_2 c,
\]
and when $a$ runs through the whole group $Q$, then, as already remarked in \S.\,1, each of the two $a a_1$ and $a a_2$ also runs through the whole group $Q$; consequently $Qb$ and $Qc$ are identical.

As in Bd.\,I, \S.\,154, 2. we can also decompose $P$ into subsidiary groups in the following way:
\[
P = Q + b^{-1} Q + c^{-1} Q + \cdots + g^{-1} Q.
\]
